\documentclass[12pt,a4paper]{article}

\usepackage[T2A]{fontenc}
\usepackage[utf8]{inputenc}
\usepackage[russian]{babel}
\usepackage{amsmath,amssymb,amsthm,mathtools}
\usepackage{enumitem}
\usepackage{array}
\usepackage{booktabs}
\usepackage{forest}
\usepackage{listings}
\usepackage{xcolor}
\usepackage{hyperref}

\hypersetup{
    colorlinks=true,
    linkcolor=blue,
    urlcolor=blue
}

\lstset{
    basicstyle=\ttfamily\small,
    keywordstyle=\color{blue},
    commentstyle=\color{gray},
    stringstyle=\color{red!60!black},
    columns=fullflexible,
    keepspaces=true,
    frame=single,
    breaklines=true
}

\newtheorem{theorem}{Теорема}
\newtheorem{lemma}{Лемма}
\newtheorem{definition}{Определение}
\newtheorem{example}{Пример}

\title{Билет №36 \\ \small Анализ исходной программы в компиляторе. Автоматные (регулярные) грамматики и сканирование, контекстно свободные грамматики и синтаксический анализ. (ИТМО) \\ Анализ исходной программы в компиляторе. Автоматные (регулярные) грамматики и сканирование, контекстно свободные грамматики и синтаксический анализ, организация таблицы символов программы, имеющей блочную структуру, хешфункции. Нисходящие (LL(1)-грамматики) и восходящие (LR(1)-грамматики) методы синтаксического анализа. Атрибутные грамматики и семантические программы, построение абстрактного синтаксического дерева. (СпБГУ)}
\author{Сериков Владислав Александрович}
\date{13 мая 2026}

\begin{document}

\maketitle
\tableofcontents
\newpage

\section{Общая структура компилятора}

\begin{definition}
\textbf{Компилятор} --- программа, переводящая исходную программу на некотором исходном языке в эквивалентную программу на целевом языке, обычно более низкого уровня.
\end{definition}

Типичная структура компилятора:

\[
\text{исходный текст}
\to
\text{лексический анализ}
\to
\text{синтаксический анализ}
\to
\text{семантический анализ}
\to
\text{промежуточное представление}
\to
\text{оптимизация}
\to
\text{генерация кода}.
\]

В рамках анализа исходной программы обычно выделяют:

\begin{enumerate}
    \item \textbf{Лексический анализ} --- преобразование потока символов в поток токенов.
    \item \textbf{Синтаксический анализ} --- проверка соответствия токенов грамматике языка и построение дерева разбора или абстрактного синтаксического дерева.
    \item \textbf{Семантический анализ} --- проверка смысловых ограничений: объявленность переменных, типы, области видимости, корректность вызовов функций и т.д.
\end{enumerate}

Пример:

\begin{lstlisting}
x = a + 10;
\end{lstlisting}

Лексический анализатор может получить токены:

\[
\langle \mathrm{id}, x\rangle,\quad
=,\quad
\langle \mathrm{id}, a\rangle,\quad
+,\quad
\langle \mathrm{num}, 10\rangle,\quad
;
\]

Синтаксический анализатор устанавливает, что это оператор присваивания, а семантический анализатор проверяет, объявлены ли $x$ и $a$, совместимы ли их типы и допустимо ли присваивание.

\section{Формальные языки и грамматики}

\subsection{Алфавит, слова, язык}

\begin{definition}
\textbf{Алфавит} $\Sigma$ --- конечное множество символов.
\end{definition}

\begin{definition}
\textbf{Слово} над алфавитом $\Sigma$ --- конечная последовательность символов из $\Sigma$.
\end{definition}

Пустое слово обозначается $\varepsilon$.

Множество всех слов над $\Sigma$ обозначается $\Sigma^*$.

\begin{definition}
\textbf{Язык} над алфавитом $\Sigma$ --- любое подмножество $L \subseteq \Sigma^*$.
\end{definition}

\subsection{Грамматика}

\begin{definition}
\textbf{Грамматика} --- четверка

\[
G = (N, \Sigma, P, S),
\]

где:
\begin{itemize}
    \item $N$ --- конечное множество нетерминалов;
    \item $\Sigma$ --- конечное множество терминалов;
    \item $P$ --- конечное множество правил вывода;
    \item $S \in N$ --- начальный нетерминал.
\end{itemize}
\end{definition}

Отношение непосредственного вывода обозначается $\Rightarrow$.

Если из $S$ можно вывести слово $w$, пишут:

\[
S \Rightarrow^* w.
\]

Язык, порождаемый грамматикой $G$:

\[
L(G) = \{w \in \Sigma^* \mid S \Rightarrow^* w\}.
\]

\section{Регулярные языки, автоматные грамматики и сканирование}

\subsection{Регулярные выражения}

Регулярные выражения описывают регулярные языки.

Базовые регулярные выражения:

\begin{itemize}
    \item $\varnothing$ --- пустой язык;
    \item $\varepsilon$ --- язык $\{\varepsilon\}$;
    \item $a$ для $a \in \Sigma$ --- язык $\{a\}$.
\end{itemize}

Операции:

\begin{itemize}
    \item объединение: $r|s$;
    \item конкатенация: $rs$;
    \item итерация Клини: $r^*$.
\end{itemize}

Пример регулярного выражения для идентификатора:

\[
\mathrm{letter}(\mathrm{letter}|\mathrm{digit})^*.
\]

\subsection{Регулярные грамматики}

\begin{definition}
\textbf{Праволинейная регулярная грамматика} --- грамматика, в которой все правила имеют один из видов:

\[
A \to aB,\qquad A \to a,\qquad A \to \varepsilon,
\]

где $A,B \in N$, $a \in \Sigma$.
\end{definition}

Аналогично определяется леволинейная регулярная грамматика:

\[
A \to Ba,\qquad A \to a,\qquad A \to \varepsilon.
\]

Регулярные грамматики используются для описания лексем: идентификаторов, чисел, ключевых слов, операторов, пробельных символов.

\subsection{Конечные автоматы}

\begin{definition}
\textbf{Детерминированный конечный автомат} --- пятерка

\[
M = (Q,\Sigma,\delta,q_0,F),
\]

где:
\begin{itemize}
    \item $Q$ --- конечное множество состояний;
    \item $\Sigma$ --- входной алфавит;
    \item $\delta: Q \times \Sigma \to Q$ --- функция переходов;
    \item $q_0 \in Q$ --- начальное состояние;
    \item $F \subseteq Q$ --- множество допускающих состояний.
\end{itemize}
\end{definition}

Автомат принимает слово $w$, если после чтения всех символов оказывается в состоянии из $F$.

\begin{definition}
\textbf{Недетерминированный конечный автомат} --- пятерка

\[
M = (Q,\Sigma,\delta,q_0,F),
\]

где

\[
\delta: Q \times (\Sigma \cup \{\varepsilon\}) \to 2^Q.
\]

То есть из одного состояния по одному символу может быть несколько переходов, а также могут существовать $\varepsilon$-переходы.
\end{definition}

\subsection{Теорема о равносильности НКА и ДКА}

\begin{theorem}
Для всякого недетерминированного конечного автомата $N$ существует детерминированный конечный автомат $D$, распознающий тот же язык:

\[
L(N)=L(D).
\]
\end{theorem}

\begin{proof}
Пусть

\[
N = (Q,\Sigma,\delta,q_0,F)
\]

--- НКА. Построим ДКА

\[
D = (Q_D,\Sigma,\delta_D,q_D,F_D).
\]

Идея состоит в том, что одно состояние ДКА будет соответствовать множеству состояний НКА, в которых НКА может находиться после чтения некоторого префикса входа.

Положим:

\[
Q_D = 2^Q.
\]

Если у НКА есть $\varepsilon$-переходы, определим $\varepsilon$-замыкание множества состояний $T$:

\[
\varepsilon\text{-closure}(T)
\]

как множество всех состояний, достижимых из состояний $T$ по нулю или более $\varepsilon$-переходам.

Начальное состояние ДКА:

\[
q_D = \varepsilon\text{-closure}(\{q_0\}).
\]

Функция переходов:

\[
\delta_D(T,a)
=
\varepsilon\text{-closure}
\left(
\bigcup_{q \in T} \delta(q,a)
\right).
\]

Допускающие состояния ДКА:

\[
F_D = \{T \subseteq Q \mid T \cap F \ne \varnothing\}.
\]

Докажем по индукции по длине слова $w$, что после чтения $w$ ДКА оказывается ровно в множестве всех состояний, в которых мог бы оказаться НКА после чтения $w$.

\textbf{База.} Для $w=\varepsilon$ ДКА находится в состоянии

\[
\varepsilon\text{-closure}(\{q_0\}),
\]

что совпадает со множеством состояний НКА, достижимых без чтения символов.

\textbf{Переход.} Пусть утверждение верно для слова $w$. Рассмотрим слово $wa$, где $a \in \Sigma$. После чтения $w$ ДКА находится в множестве $T$, равном множеству всех возможных состояний НКА после $w$. После чтения символа $a$ НКА может перейти в любое состояние из

\[
\bigcup_{q \in T}\delta(q,a),
\]

а затем пройти по $\varepsilon$-переходам. Именно это и задает

\[
\delta_D(T,a).
\]

Следовательно, индукционное утверждение верно для $wa$.

Таким образом, после чтения любого слова $w$ ДКА хранит множество возможных состояний НКА после чтения $w$. Слово принимается НКА тогда и только тогда, когда среди этих состояний есть хотя бы одно допускающее состояние, то есть когда соответствующее состояние ДКА принадлежит $F_D$.

Следовательно,

\[
L(N)=L(D).
\]
\end{proof}

\subsection{Теорема Клини}

\begin{theorem}[Теорема Клини]
Класс языков, задаваемых регулярными выражениями, совпадает с классом языков, распознаваемых конечными автоматами.
\end{theorem}

\begin{proof}
Доказательство состоит из двух частей.

\textbf{1. Из регулярного выражения можно построить конечный автомат.}

Покажем это структурной индукцией по регулярному выражению.

Для базовых выражений:
\begin{itemize}
    \item для $\varnothing$ строится автомат без принимающих путей;
    \item для $\varepsilon$ строится автомат, принимающий пустую строку;
    \item для символа $a$ строится автомат с переходом из начального состояния в допускающее по символу $a$.
\end{itemize}

Далее пусть уже построены автоматы для регулярных выражений $r$ и $s$.

\begin{itemize}
    \item Для $r|s$ создается новое начальное состояние, из которого $\varepsilon$-переходами можно попасть в начальные состояния автоматов для $r$ и $s$.
    \item Для $rs$ допускающие состояния автомата для $r$ соединяются $\varepsilon$-переходами с начальным состоянием автомата для $s$.
    \item Для $r^*$ добавляются новое начальное и новое допускающее состояния; из начального состояния есть $\varepsilon$-переход в автомат для $r$ и в новое допускающее состояние, а из допускающих состояний автомата для $r$ есть $\varepsilon$-переходы обратно к началу автомата для $r$ и к новому допускающему состоянию.
\end{itemize}

Следовательно, для любого регулярного выражения можно построить НКА, принимающий тот же язык. По теореме о детерминизации НКА существует эквивалентный ДКА.

\textbf{2. Из конечного автомата можно построить регулярное выражение.}

Пусть дан конечный автомат. Построим обобщенный конечный автомат, в котором переходы помечены регулярными выражениями. Затем будем последовательно устранять состояния.

Если устраняется состояние $k$, то для каждой пары состояний $i,j$ обновляем метку перехода:

\[
R_{ij} \gets R_{ij} \mid R_{ik}(R_{kk})^*R_{kj}.
\]

Смысл формулы: путь из $i$ в $j$ либо не проходит через $k$, либо проходит через $k$ один или несколько раз, причем циклы в $k$ описываются $(R_{kk})^*$.

После устранения всех промежуточных состояний остается регулярное выражение на переходе из начального состояния в допускающее, описывающее тот же язык.

Итак, каждый регулярный язык распознается конечным автоматом, и каждый язык конечного автомата задается регулярным выражением. Следовательно, классы совпадают.
\end{proof}

\subsection{Лексический анализ}

\begin{definition}
\textbf{Лексический анализатор}, или \textbf{сканер}, преобразует поток символов исходной программы в поток токенов.
\end{definition}

\begin{definition}
\textbf{Токен} --- класс лексем, например \texttt{IDENT}, \texttt{NUMBER}, \texttt{IF}, \texttt{PLUS}.
\end{definition}

\begin{definition}
\textbf{Лексема} --- конкретная подстрока исходного текста, соответствующая токену.
\end{definition}

Пример:

\begin{lstlisting}
if (x1 >= 10) x1 = x1 + 1;
\end{lstlisting}

Токены:

\[
\mathrm{IF},\quad
(,\quad
\langle \mathrm{IDENT},x1\rangle,\quad
\ge,\quad
\langle \mathrm{NUMBER},10\rangle,\quad
),\quad
\langle \mathrm{IDENT},x1\rangle,\quad
=,\quad
\langle \mathrm{IDENT},x1\rangle,\quad
+,\quad
\langle \mathrm{NUMBER},1\rangle,\quad
;.
\]

\subsection{Типичный алгоритм сканирования}

\textbf{Вход:} строка исходной программы.

\textbf{Выход:} последовательность токенов.

\begin{enumerate}
    \item Установить указатель \texttt{pos} на начало входной строки.
    \item Пока \texttt{pos} не достиг конца строки:
    \begin{enumerate}
        \item Пропустить пробелы, табуляции, переводы строк, если они не являются значимыми.
        \item Из текущей позиции запустить конечный автомат для распознавания токенов.
        \item Прочитать максимальную подстроку, соответствующую некоторому токену.
        \item Если несколько правил подходят, выбрать правило с наибольшим приоритетом.
        \item Вернуть соответствующий токен.
        \item Передвинуть \texttt{pos} за конец найденной лексемы.
    \end{enumerate}
    \item Вернуть специальный токен конца файла \texttt{EOF}.
\end{enumerate}

Два стандартных правила сканирования:

\begin{enumerate}
    \item \textbf{Правило максимального соответствия}:
    выбирается самая длинная лексема.
    \item \textbf{Правило приоритета}:
    если несколько токенов имеют одинаковую длину, выбирается токен, правило которого указано раньше.
\end{enumerate}

Пример:

\begin{lstlisting}
ifx
\end{lstlisting}

Если ключевое слово \texttt{if} и идентификатор задаются правилами:

\[
\mathrm{IF}: \texttt{if},
\qquad
\mathrm{IDENT}: [a-zA-Z][a-zA-Z0-9]^*,
\]

то строка \texttt{ifx} распознается как один идентификатор, а не как \texttt{IF IDENT}.

\section{Контекстно-свободные грамматики и синтаксический анализ}

\subsection{Контекстно-свободные грамматики}

\begin{definition}
\textbf{Контекстно-свободная грамматика} --- грамматика

\[
G=(N,\Sigma,P,S),
\]

в которой каждое правило имеет вид

\[
A \to \alpha,
\]

где $A \in N$, $\alpha \in (N \cup \Sigma)^*$.
\end{definition}

Контекстно-свободные грамматики описывают вложенные структуры: выражения, блоки, условные операторы, циклы, списки параметров.

Пример грамматики арифметических выражений:

\[
\begin{aligned}
E &\to E + T \mid T,\\
T &\to T * F \mid F,\\
F &\to (E) \mid \mathrm{id}.
\end{aligned}
\]

Эта грамматика задает приоритет умножения над сложением.

\subsection{Выводы}

\begin{definition}
\textbf{Левосторонний вывод} --- вывод, в котором на каждом шаге заменяется самый левый нетерминал.
\end{definition}

\begin{definition}
\textbf{Правосторонний вывод} --- вывод, в котором на каждом шаге заменяется самый правый нетерминал.
\end{definition}

Для грамматики:

\[
E \to E+T \mid T,\quad T \to \mathrm{id}
\]

левосторонний вывод строки $\mathrm{id}+\mathrm{id}$:

\[
E \Rightarrow E+T \Rightarrow T+T \Rightarrow \mathrm{id}+T \Rightarrow \mathrm{id}+\mathrm{id}.
\]

\subsection{Дерево разбора}

\begin{definition}
\textbf{Дерево разбора} --- дерево, в котором:
\begin{itemize}
    \item корень помечен начальным символом;
    \item внутренние вершины помечены нетерминалами;
    \item листья помечены терминалами или $\varepsilon$;
    \item дети каждой внутренней вершины соответствуют правой части примененного правила.
\end{itemize}
\end{definition}

Пример для выражения:

\[
\mathrm{id}+\mathrm{id}*\mathrm{id}.
\]

\[
\begin{forest}
[E
  [E
    [T
      [F [id]]
    ]
  ]
  [+]
  [T
    [T
      [F [id]]
    ]
    [*]
    [F [id]]
  ]
]
\end{forest}
\]

\subsection{Неоднозначность грамматики}

\begin{definition}
Контекстно-свободная грамматика называется \textbf{неоднозначной}, если существует хотя бы одна строка, имеющая два различных дерева разбора.
\end{definition}

Пример неоднозначной грамматики:

\[
E \to E+E \mid E*E \mid (E) \mid \mathrm{id}.
\]

Строка

\[
\mathrm{id}+\mathrm{id}*\mathrm{id}
\]

имеет два дерева разбора:

\[
(\mathrm{id}+\mathrm{id})*\mathrm{id}
\]

и

\[
\mathrm{id}+(\mathrm{id}*\mathrm{id}).
\]

Для компиляторов неоднозначность нежелательна, поскольку синтаксическая структура должна быть определена однозначно.

\section{FIRST и FOLLOW}

\subsection{Множество FIRST}

\begin{definition}
Для строки грамматических символов $\alpha$ множество $\mathrm{FIRST}(\alpha)$ --- множество терминалов, с которых могут начинаться строки, выводимые из $\alpha$:

\[
\mathrm{FIRST}(\alpha)
=
\{a \in \Sigma \mid \alpha \Rightarrow^* a\beta\}.
\]

Если $\alpha \Rightarrow^* \varepsilon$, то $\varepsilon \in \mathrm{FIRST}(\alpha)$.
\end{definition}

\subsection{Множество FOLLOW}

\begin{definition}
Для нетерминала $A$ множество $\mathrm{FOLLOW}(A)$ --- множество терминалов, которые могут непосредственно следовать за $A$ в некоторой сентенциальной форме:

\[
\mathrm{FOLLOW}(A)
=
\{a \in \Sigma \mid S \Rightarrow^* \alpha A a \beta\}.
\]

Для начального символа $S$ в $\mathrm{FOLLOW}(S)$ включают маркер конца ввода $\$$.
\end{definition}

\subsection{Алгоритм вычисления FIRST}

\textbf{Вход:} КС-грамматика $G$.

\textbf{Выход:} множества $\mathrm{FIRST}$ для всех символов.

\begin{enumerate}
    \item Для каждого терминала $a$ положить:
    \[
    \mathrm{FIRST}(a)=\{a\}.
    \]
    \item Для каждого нетерминала $A$ положить:
    \[
    \mathrm{FIRST}(A)=\varnothing.
    \]
    \item Повторять, пока множества изменяются:
    \begin{enumerate}
        \item Для каждого правила
        \[
        A \to X_1X_2\ldots X_k
        \]
        добавить в $\mathrm{FIRST}(A)$ все терминалы из $\mathrm{FIRST}(X_1)$, кроме $\varepsilon$.
        \item Если $\varepsilon \in \mathrm{FIRST}(X_1)$, добавить терминалы из $\mathrm{FIRST}(X_2)$, кроме $\varepsilon$, и так далее.
        \item Если для всех $i$ выполнено $\varepsilon \in \mathrm{FIRST}(X_i)$, добавить $\varepsilon$ в $\mathrm{FIRST}(A)$.
    \end{enumerate}
\end{enumerate}

\subsection{Алгоритм вычисления FOLLOW}

\begin{enumerate}
    \item Для всех нетерминалов $A$ положить:
    \[
    \mathrm{FOLLOW}(A)=\varnothing.
    \]
    \item Добавить $\$$ в $\mathrm{FOLLOW}(S)$.
    \item Повторять, пока множества изменяются:
    \begin{enumerate}
        \item Для каждого правила
        \[
        A \to \alpha B \beta
        \]
        добавить
        \[
        \mathrm{FIRST}(\beta)\setminus\{\varepsilon\}
        \]
        в $\mathrm{FOLLOW}(B)$.
        \item Если $\varepsilon \in \mathrm{FIRST}(\beta)$ или $\beta=\varepsilon$, добавить
        \[
        \mathrm{FOLLOW}(A)
        \]
        в $\mathrm{FOLLOW}(B)$.
    \end{enumerate}
\end{enumerate}

\subsection{Минимальный пример}

Пусть дана грамматика:

\[
\begin{aligned}
S &\to AB,\\
A &\to a \mid \varepsilon,\\
B &\to b.
\end{aligned}
\]

Тогда:

\[
\mathrm{FIRST}(A)=\{a,\varepsilon\},
\qquad
\mathrm{FIRST}(B)=\{b\}.
\]

Поскольку $S \to AB$, имеем:

\[
\mathrm{FIRST}(S)=\{a,b\}.
\]

Множества FOLLOW:

\[
\mathrm{FOLLOW}(S)=\{\$\}.
\]

Так как в правиле $S \to AB$ после $A$ стоит $B$, то:

\[
\mathrm{FOLLOW}(A)=\{b\}.
\]

Так как $B$ стоит в конце правила $S \to AB$, то:

\[
\mathrm{FOLLOW}(B)=\{\$\}.
\]

\section{Нисходящий синтаксический анализ и LL(1)-грамматики}

\subsection{Идея нисходящего анализа}

Нисходящий синтаксический анализ строит дерево разбора от корня к листьям. Он пытается вывести входную строку из начального символа.

Наиболее важный практический класс --- LL(1)-грамматики.

\begin{definition}
Грамматика называется \textbf{LL(1)-грамматикой}, если:
\begin{itemize}
    \item первая буква L означает чтение входа слева направо;
    \item вторая буква L означает построение левостороннего вывода;
    \item число 1 означает использование одного символа предпросмотра.
\end{itemize}
\end{definition}

\subsection{Условие LL(1)}

Для каждого нетерминала $A$ и любых двух различных продукций:

\[
A \to \alpha,
\qquad
A \to \beta
\]

должны выполняться условия:

\[
\mathrm{FIRST}(\alpha)\cap \mathrm{FIRST}(\beta)=\varnothing.
\]

Если $\varepsilon \in \mathrm{FIRST}(\alpha)$, то дополнительно:

\[
\mathrm{FIRST}(\beta)\cap \mathrm{FOLLOW}(A)=\varnothing.
\]

Если $\varepsilon \in \mathrm{FIRST}(\beta)$, то:

\[
\mathrm{FIRST}(\alpha)\cap \mathrm{FOLLOW}(A)=\varnothing.
\]

\subsection{Теорема о критерии LL(1)}

\begin{theorem}
Контекстно-свободная грамматика является LL(1)-грамматикой тогда и только тогда, когда для любых двух различных правил одного нетерминала

\[
A \to \alpha,
\qquad
A \to \beta
\]

выполнены условия:
\[
\mathrm{FIRST}(\alpha)\cap \mathrm{FIRST}(\beta)=\varnothing,
\]
а если одна из правых частей выводит $\varepsilon$, то множество FIRST другой правой части не пересекается с $\mathrm{FOLLOW}(A)$.
\end{theorem}

\begin{proof}
Докажем необходимость и достаточность.

\textbf{Необходимость.}
Пусть грамматика является LL(1). Тогда при разборе нетерминала $A$ и одном символе предпросмотра парсер обязан однозначно выбрать одно правило.

Если существует терминал

\[
a \in \mathrm{FIRST}(\alpha)\cap \mathrm{FIRST}(\beta),
\]

то обе правые части могут породить строки, начинающиеся с $a$. Тогда по символу предпросмотра $a$ парсер не сможет выбрать между правилами $A \to \alpha$ и $A \to \beta$. Противоречие с LL(1).

Теперь пусть $\varepsilon \in \mathrm{FIRST}(\alpha)$ и существует

\[
a \in \mathrm{FIRST}(\beta)\cap \mathrm{FOLLOW}(A).
\]

Тогда при символе предпросмотра $a$ возможны два варианта:
\begin{itemize}
    \item применить $A \to \alpha$, вывести $\varepsilon$ и позволить символу $a$ следовать за $A$;
    \item применить $A \to \beta$, поскольку $\beta$ может породить строку, начинающуюся с $a$.
\end{itemize}

Снова возникает неоднозначность выбора. Следовательно, пересечение должно быть пустым.

\textbf{Достаточность.}
Пусть указанные условия выполнены. При разборе нетерминала $A$ с символом предпросмотра $a$ парсер выбирает правило $A \to \gamma$ так:
\begin{itemize}
    \item если $a \in \mathrm{FIRST}(\gamma)$, выбирается $A \to \gamma$;
    \item если $\varepsilon \in \mathrm{FIRST}(\gamma)$ и $a \in \mathrm{FOLLOW}(A)$, выбирается $A \to \gamma$.
\end{itemize}

По условиям теоремы никакое другое правило для $A$ не может быть применимо к тому же символу предпросмотра. Значит, в каждой ячейке LL(1)-таблицы находится не более одного правила. Следовательно, разбор детерминирован при одном символе предпросмотра, то есть грамматика является LL(1).
\end{proof}

\subsection{Построение LL(1)-таблицы}

\textbf{Вход:} грамматика $G$.

\textbf{Выход:} таблица разбора $M[A,a]$.

\begin{enumerate}
    \item Вычислить множества $\mathrm{FIRST}$ и $\mathrm{FOLLOW}$.
    \item Для каждого правила
    \[
    A \to \alpha
    \]
    выполнить:
    \begin{enumerate}
        \item Для каждого терминала
        \[
        a \in \mathrm{FIRST}(\alpha)\setminus\{\varepsilon\}
        \]
        записать правило $A \to \alpha$ в ячейку $M[A,a]$.
        \item Если
        \[
        \varepsilon \in \mathrm{FIRST}(\alpha),
        \]
        то для каждого
        \[
        b \in \mathrm{FOLLOW}(A)
        \]
        записать правило $A \to \alpha$ в $M[A,b]$.
    \end{enumerate}
    \item Если хотя бы в одной ячейке оказалось более одного правила, грамматика не является LL(1).
\end{enumerate}

\subsection{Пример LL(1)-грамматики}

Возьмем грамматику:

\[
\begin{aligned}
E &\to TE',\\
E' &\to +TE' \mid \varepsilon,\\
T &\to id.
\end{aligned}
\]

Множества:

\[
\mathrm{FIRST}(E)=\{id\},
\quad
\mathrm{FIRST}(E')=\{+,\varepsilon\},
\quad
\mathrm{FIRST}(T)=\{id\}.
\]

\[
\mathrm{FOLLOW}(E)=\{\$,)\},
\quad
\mathrm{FOLLOW}(E')=\{\$,)\},
\quad
\mathrm{FOLLOW}(T)=\{+,\$,)\}.
\]

Фрагмент таблицы:

\[
\begin{array}{c|ccc}
 & id & + & \$ \\
\hline
E  & E\to TE' & & \\
E' & & E'\to +TE' & E'\to \varepsilon \\
T  & T\to id & & \\
\end{array}
\]

\subsection{Алгоритм табличного LL(1)-разбора}

\textbf{Вход:} строка токенов $w\$$.

\textbf{Стек:} изначально содержит $\$S$.

\begin{enumerate}
    \item Пусть $X$ --- верхушка стека, $a$ --- текущий входной символ.
    \item Если $X$ --- терминал и $X=a$, снять $X$ со стека и сдвинуть вход.
    \item Если $X$ --- терминал и $X\ne a$, сообщить об ошибке.
    \item Если $X$ --- нетерминал, посмотреть таблицу $M[X,a]$.
    \begin{enumerate}
        \item Если ячейка пуста, сообщить об ошибке.
        \item Если в ячейке правило $X \to Y_1Y_2\ldots Y_k$, снять $X$ и положить $Y_k,\ldots,Y_2,Y_1$ на стек.
        \item Если правило $X \to \varepsilon$, просто снять $X$.
    \end{enumerate}
    \item Если $X=\$$ и $a=\$$, принять строку.
\end{enumerate}

\section{Преобразования грамматик для нисходящего анализа}

\subsection{Устранение левой рекурсии}

\begin{definition}
Грамматика имеет \textbf{непосредственную левую рекурсию}, если есть правило вида:

\[
A \to A\alpha.
\]
\end{definition}

Левая рекурсия мешает рекурсивному спуску, потому что процедура для $A$ вызывает саму себя без потребления входа.

Если есть правила:

\[
A \to A\alpha_1 \mid A\alpha_2 \mid \ldots \mid A\alpha_m
\mid
\beta_1 \mid \beta_2 \mid \ldots \mid \beta_n,
\]

где $\beta_i$ не начинаются с $A$, то заменяем их на:

\[
A \to \beta_1 A' \mid \beta_2 A' \mid \ldots \mid \beta_n A',
\]

\[
A' \to \alpha_1 A' \mid \alpha_2 A' \mid \ldots \mid \alpha_m A' \mid \varepsilon.
\]

Пример:

\[
E \to E+T \mid T
\]

заменяется на:

\[
E \to TE',
\qquad
E' \to +TE' \mid \varepsilon.
\]

\subsection{Левая факторизация}

Если у нескольких правил есть общий префикс:

\[
A \to \alpha\beta_1 \mid \alpha\beta_2,
\]

то парсер с одним символом предпросмотра может не понять, какое правило выбрать.

Преобразование:

\[
A \to \alpha A',
\qquad
A' \to \beta_1 \mid \beta_2.
\]

Пример:

\[
S \to if\ E\ then\ S\ else\ S
\mid
if\ E\ then\ S
\]

после факторизации:

\[
S \to if\ E\ then\ S\ S',
\]

\[
S' \to else\ S \mid \varepsilon.
\]

\section{Восходящий синтаксический анализ и LR(1)-грамматики}

\subsection{Идея восходящего анализа}

Восходящий анализ строит дерево разбора от листьев к корню. Он пытается свернуть входную строку к начальному символу грамматики.

Основная операция --- \textbf{сдвиг-свертка}.

\begin{itemize}
    \item \textbf{Сдвиг} --- перенос очередного входного символа в стек.
    \item \textbf{Свертка} --- замена правой части некоторого правила на его левую часть.
\end{itemize}

Пример:

\[
E \to E+T \mid T,\qquad T\to id.
\]

Строка:

\[
id+id.
\]

Возможные шаги:

\[
id+id
\Rightarrow T+id
\Rightarrow E+id
\Rightarrow E+T
\Rightarrow E.
\]

В терминах стека:

\[
\begin{array}{c|c|c}
\text{Стек} & \text{Вход} & \text{Действие} \\
\hline
\varepsilon & id+id\$ & shift \\
id & +id\$ & reduce\ T\to id \\
T & +id\$ & reduce\ E\to T \\
E & +id\$ & shift \\
E+ & id\$ & shift \\
E+id & \$ & reduce\ T\to id \\
E+T & \$ & reduce\ E\to E+T \\
E & \$ & accept \\
\end{array}
\]

\subsection{LR(1)-грамматики}

\begin{definition}
Грамматика называется \textbf{LR(1)-грамматикой}, если:
\begin{itemize}
    \item L означает чтение входа слева направо;
    \item R означает построение правостороннего вывода в обратном порядке;
    \item 1 означает использование одного символа предпросмотра.
\end{itemize}
\end{definition}

LR-анализаторы более мощные, чем LL-анализаторы. Они способны обрабатывать многие грамматики с левой рекурсией, что удобно для выражений.

\subsection{LR(1)-пункт}

\begin{definition}
\textbf{LR(1)-пункт} --- запись вида

\[
[A \to \alpha \cdot \beta, a],
\]

где:
\begin{itemize}
    \item $A \to \alpha\beta$ --- правило грамматики;
    \item точка показывает, какая часть правой части уже распознана;
    \item $a$ --- символ предпросмотра.
\end{itemize}
\end{definition}

Например:

\[
[E \to E + \cdot T, \$]
\]

означает, что уже распознано $E+$, дальше ожидается $T$, а после завершения $E$ ожидается конец ввода.

\subsection{Операция CLOSURE}

Пусть $I$ --- множество LR(1)-пунктов.

\[
\mathrm{CLOSURE}(I)
\]

строится так:

\begin{enumerate}
    \item Изначально положить $\mathrm{CLOSURE}(I)=I$.
    \item Если в множестве есть пункт
    \[
    [A \to \alpha \cdot B\beta, a],
    \]
    где после точки стоит нетерминал $B$, то для каждого правила
    \[
    B \to \gamma
    \]
    и каждого терминала
    \[
    b \in \mathrm{FIRST}(\beta a)
    \]
    добавить пункт
    \[
    [B \to \cdot \gamma, b].
    \]
    \item Повторять, пока новые пункты больше не добавляются.
\end{enumerate}

\subsection{Операция GOTO}

Для множества пунктов $I$ и грамматического символа $X$:

\[
\mathrm{GOTO}(I,X)
=
\mathrm{CLOSURE}
\left(
\{[A\to \alpha X \cdot \beta,a]
\mid
[A\to \alpha \cdot X\beta,a]\in I\}
\right).
\]

Иначе говоря, если в состоянии можно прочитать символ $X$, точка переносится через $X$, а затем множество пунктов замыкается.

\subsection{Построение канонического множества LR(1)-состояний}

\begin{enumerate}
    \item Добавить новое начальное правило:
    \[
    S' \to S.
    \]
    \item Построить начальное состояние:
    \[
    I_0 = \mathrm{CLOSURE}(\{[S'\to \cdot S,\$]\}).
    \]
    \item Для каждого уже построенного состояния $I$ и каждого грамматического символа $X$ вычислить:
    \[
    J=\mathrm{GOTO}(I,X).
    \]
    \item Если $J\ne\varnothing$ и такого состояния еще нет, добавить его.
    \item Повторять до неподвижной точки.
\end{enumerate}

\subsection{Построение таблиц ACTION и GOTO}

LR-анализатор использует две таблицы:

\begin{itemize}
    \item \texttt{ACTION[state, terminal]} --- действие: shift, reduce, accept или error;
    \item \texttt{GOTO[state, nonterminal]} --- переход по нетерминалу после свертки.
\end{itemize}

Правила заполнения:

\begin{enumerate}
    \item Если в состоянии $I_i$ есть пункт
    \[
    [A \to \alpha \cdot a\beta, b]
    \]
    и
    \[
    \mathrm{GOTO}(I_i,a)=I_j,
    \]
    где $a$ --- терминал, то:
    \[
    \mathrm{ACTION}[i,a]=shift\ j.
    \]
    \item Если в состоянии $I_i$ есть пункт
    \[
    [A \to \alpha \cdot, a],
    \]
    где $A\ne S'$, то:
    \[
    \mathrm{ACTION}[i,a]=reduce\ A\to \alpha.
    \]
    \item Если в состоянии $I_i$ есть пункт
    \[
    [S'\to S\cdot,\$],
    \]
    то:
    \[
    \mathrm{ACTION}[i,\$]=accept.
    \]
    \item Если
    \[
    \mathrm{GOTO}(I_i,A)=I_j
    \]
    для нетерминала $A$, то:
    \[
    \mathrm{GOTO}[i,A]=j.
    \]
\end{enumerate}

Если при заполнении таблицы в одной ячейке возникает несколько действий, появляется конфликт:

\begin{itemize}
    \item shift/reduce;
    \item reduce/reduce.
\end{itemize}

Грамматика является LR(1), если ее каноническая LR(1)-таблица не содержит конфликтов.

\subsection{Алгоритм LR-разбора}

\textbf{Стек} содержит пары вида состояние/символ. Обычно хранят состояния.

\begin{enumerate}
    \item Изначально стек содержит состояние $0$.
    \item Пусть $s$ --- верхнее состояние стека, $a$ --- текущий входной символ.
    \item Посмотреть \texttt{ACTION[$s$,$a$]}.
    \item Если действие \texttt{shift $t$}:
    \begin{enumerate}
        \item поместить символ $a$ и состояние $t$ в стек;
        \item сдвинуть вход.
    \end{enumerate}
    \item Если действие \texttt{reduce $A\to \beta$}:
    \begin{enumerate}
        \item снять со стека $|\beta|$ грамматических символов и соответствующие состояния;
        \item пусть $s'$ --- новое верхнее состояние;
        \item положить $A$;
        \item перейти в состояние
        \[
        \mathrm{GOTO}[s',A].
        \]
    \end{enumerate}
    \item Если действие \texttt{accept}, принять строку.
    \item Если действие отсутствует, сообщить об ошибке.
\end{enumerate}

\section{Сравнение LL(1) и LR(1)}

\[
\begin{array}{p{0.28\textwidth}|p{0.32\textwidth}|p{0.32\textwidth}}
\toprule
\textbf{Свойство} & \textbf{LL(1)} & \textbf{LR(1)} \\
\midrule
Направление построения дерева & Сверху вниз & Снизу вверх \\
Вывод & Левосторонний & Правосторонний в обратном порядке \\
Левая рекурсия & Недопустима & Допустима \\
Мощность & Менее мощный класс & Более мощный класс \\
Реализация вручную & Удобна & Сложнее \\
Таблицы & Обычно меньше и проще & Обычно больше \\
Типичные ошибки & Раннее обнаружение & Хорошая локализация в практических реализациях \\
\bottomrule
\end{array}
\]

\section{Атрибутные грамматики и семантические программы}

\subsection{Атрибутная грамматика}

\begin{definition}
\textbf{Атрибутная грамматика} --- контекстно-свободная грамматика, символам которой сопоставлены атрибуты, а правилам --- семантические правила для вычисления этих атрибутов.
\end{definition}

Атрибуты бывают:

\begin{itemize}
    \item \textbf{синтезируемые} --- вычисляются от дочерних узлов к родительскому;
    \item \textbf{наследуемые} --- передаются от родителя или соседних узлов к дочернему.
\end{itemize}

Пример синтезируемого атрибута для вычисления значения выражения:

\[
E \to E_1 + T
\]

Семантическое правило:

\[
E.val = E_1.val + T.val.
\]

Для правила:

\[
T \to \mathrm{num}
\]

\[
T.val = \mathrm{num.lexval}.
\]

\subsection{S-атрибутные грамматики}

\begin{definition}
\textbf{S-атрибутная грамматика} --- атрибутная грамматика, в которой все атрибуты синтезируемые.
\end{definition}

Такие грамматики удобно вычислять при восходящем анализе: когда выполняется свертка, значения атрибутов правой части уже известны.

Пример:

\[
\begin{aligned}
E &\to E_1 + T
    && E.val = E_1.val + T.val,\\
E &\to T
    && E.val = T.val,\\
T &\to \mathrm{num}
    && T.val = \mathrm{num.lexval}.
\end{aligned}
\]

\subsection{L-атрибутные грамматики}

\begin{definition}
Атрибутная грамматика называется \textbf{L-атрибутной}, если каждый наследуемый атрибут символа $X_i$ в правой части правила

\[
A \to X_1X_2\ldots X_n
\]

зависит только от:
\begin{itemize}
    \item наследуемых атрибутов $A$;
    \item атрибутов символов $X_1,\ldots,X_{i-1}$, стоящих левее $X_i$.
\end{itemize}
\end{definition}

L-атрибутные грамматики удобно вычислять при нисходящем анализе слева направо.

\subsection{Синтаксически управляемый перевод}

\begin{definition}
\textbf{Синтаксически управляемый перевод} --- метод, при котором действия семантического анализа и генерации промежуточного кода привязываются к правилам грамматики.
\end{definition}

Пример генерации постфиксной записи:

\[
E \to E_1 + T
\]

Семантическое действие:

\[
\texttt{print("+")}.
\]

Для выражения:

\[
a+b*c
\]

результат:

\[
a\ b\ c\ *\ +.
\]

\section{Абстрактное синтаксическое дерево}

\subsection{Дерево разбора и AST}

Дерево разбора отражает все детали грамматики: нетерминалы, скобки, вспомогательные символы.

\begin{definition}
\textbf{Абстрактное синтаксическое дерево} --- дерево, представляющее существенную синтаксическую структуру программы без лишних деталей конкретной грамматики.
\end{definition}

Например, выражение:

\[
a + b * c
\]

может иметь AST:

\[
\begin{forest}
[+
  [a]
  [*
    [b]
    [c]
  ]
]
\end{forest}
\]

Скобки и промежуточные нетерминалы не сохраняются, если они не нужны для семантики.

\subsection{Построение AST с помощью семантических действий}

Для грамматики:

\[
\begin{aligned}
E &\to E_1 + T,\\
E &\to T,\\
T &\to T_1 * F,\\
T &\to F,\\
F &\to id.
\end{aligned}
\]

Можно задать действия:

\[
E.node = \mathrm{Node}("+", E_1.node, T.node),
\]

\[
T.node = \mathrm{Node}("*", T_1.node, F.node),
\]

\[
F.node = \mathrm{Leaf}(id.lexeme).
\]

Минимальный пример для выражения:

\[
a+b*c.
\]

Шаги построения:

\begin{enumerate}
    \item Для \texttt{a}, \texttt{b}, \texttt{c} создаются листья:
    \[
    \mathrm{Leaf}(a),\quad \mathrm{Leaf}(b),\quad \mathrm{Leaf}(c).
    \]
    \item Для \texttt{b*c} создается узел:
    \[
    \mathrm{Node}(*,\mathrm{Leaf}(b),\mathrm{Leaf}(c)).
    \]
    \item Для всего выражения создается узел:
    \[
    \mathrm{Node}(+,\mathrm{Leaf}(a),\mathrm{Node}(*,\mathrm{Leaf}(b),\mathrm{Leaf}(c))).
    \]
\end{enumerate}

\section{Семантический анализ}

Синтаксис не выражает все ограничения языка. Например, строка:

\begin{lstlisting}
x = y + 1;
\end{lstlisting}

может быть синтаксически корректной, но семантически некорректной, если переменная \texttt{y} не объявлена.

Основные задачи семантического анализа:

\begin{enumerate}
    \item Проверка объявленности идентификаторов.
    \item Проверка областей видимости.
    \item Проверка типов.
    \item Проверка числа и типов аргументов функций.
    \item Проверка допустимости операций.
    \item Разрешение перегрузки.
    \item Построение и использование таблицы символов.
\end{enumerate}

\section{Таблица символов}

\subsection{Назначение таблицы символов}

\begin{definition}
\textbf{Таблица символов} --- структура данных, содержащая информацию об идентификаторах программы.
\end{definition}

Для каждого имени могут храниться:

\begin{itemize}
    \item имя;
    \item категория: переменная, функция, тип, поле класса;
    \item тип;
    \item область видимости;
    \item адрес или смещение;
    \item параметры функции;
    \item модификаторы;
    \item ссылка на узел AST.
\end{itemize}

Пример записи:

\[
\begin{array}{c|c|c|c}
\text{Имя} & \text{Категория} & \text{Тип} & \text{Смещение} \\
\hline
x & variable & int & -4 \\
f & function & int(int,int) & label_f \\
\end{array}
\]

\subsection{Блочная структура программы}

Во многих языках блоки образуют вложенные области видимости.

Пример:

\begin{lstlisting}
int x;

void f() {
    int y;
    {
        int x;
        y = x;
    }
}
\end{lstlisting}

Внутреннее объявление \texttt{x} скрывает внешнее.

Области видимости можно представить деревом:

\[
\begin{forest}
[global
  [function f
    [inner block]
  ]
]
\end{forest}
\]

\subsection{Организация таблицы символов для блоков}

Распространенный способ --- стек таблиц символов.

\begin{enumerate}
    \item При входе в новый блок создается новая таблица и помещается на стек.
    \item При объявлении имени оно вставляется в верхнюю таблицу.
    \item При использовании имени поиск идет сверху вниз:
    \[
    \text{текущий блок} \to \text{родительский блок} \to \ldots \to \text{глобальная область}.
    \]
    \item При выходе из блока верхняя таблица снимается со стека.
\end{enumerate}

\subsection{Алгоритм вставки имени}

\textbf{Вход:} имя $x$, информация о символе.

\begin{enumerate}
    \item Взять верхнюю таблицу стека областей видимости.
    \item Проверить, есть ли в ней имя $x$.
    \item Если есть, сообщить об ошибке повторного объявления в одной области.
    \item Если нет, добавить запись для $x$.
\end{enumerate}

\subsection{Алгоритм поиска имени}

\textbf{Вход:} имя $x$.

\textbf{Выход:} запись таблицы символов или ошибка.

\begin{enumerate}
    \item Начать с верхней таблицы стека.
    \item Если имя $x$ найдено, вернуть соответствующую запись.
    \item Иначе перейти к следующей таблице ниже в стеке.
    \item Если достигнут конец стека, сообщить об ошибке: имя не объявлено.
\end{enumerate}

Пример:

\begin{lstlisting}
int x;

{
    int y;
    {
        int x;
        x = y;
    }
}
\end{lstlisting}

При обработке \texttt{x = y}:

\begin{itemize}
    \item поиск \texttt{x} найдет внутреннее объявление;
    \item поиск \texttt{y} сначала не найдет имя во внутреннем блоке, затем найдет его во внешнем блоке.
\end{itemize}

\section{Хеш-таблицы и хеш-функции}

\subsection{Хеш-таблица}

\begin{definition}
\textbf{Хеш-таблица} --- структура данных, реализующая ассоциативный массив и поддерживающая операции вставки, поиска и удаления в среднем за $O(1)$.
\end{definition}

Идея:

\[
h(k) \to \{0,1,\ldots,m-1\},
\]

где:
\begin{itemize}
    \item $k$ --- ключ, например имя идентификатора;
    \item $m$ --- размер таблицы;
    \item $h$ --- хеш-функция.
\end{itemize}

\subsection{Хеш-функция}

Для строки $s=s_0s_1\ldots s_{n-1}$ часто используют полиномиальное хеширование:

\[
h(s)
=
\left(
\sum_{i=0}^{n-1}
s_i p^i
\right)
\bmod m,
\]

где $p$ --- основание, $m$ --- размер таблицы.

Практический вариант:

\[
h = 0,
\]

\[
h \gets (h \cdot p + \mathrm{code}(c)) \bmod m
\]

для каждого символа $c$.

\subsection{Коллизии}

\begin{definition}
\textbf{Коллизия} --- ситуация, когда два разных ключа имеют одинаковое значение хеш-функции.
\end{definition}

Основные методы разрешения коллизий:

\begin{enumerate}
    \item Метод цепочек.
    \item Открытая адресация.
\end{enumerate}

\subsection{Метод цепочек}

Каждая ячейка таблицы содержит список элементов.

\[
\begin{array}{c|l}
0 & \\
1 & (abc) \to (cab) \\
2 & \\
3 & (foo) \\
\end{array}
\]

Вставка:

\begin{enumerate}
    \item Вычислить $i=h(k)$.
    \item Просмотреть список в ячейке $i$.
    \item Если ключ уже есть, обновить или сообщить об ошибке.
    \item Иначе добавить элемент в список.
\end{enumerate}

Поиск:

\begin{enumerate}
    \item Вычислить $i=h(k)$.
    \item Просмотреть список в ячейке $i$.
    \item Если ключ найден, вернуть значение.
    \item Иначе сообщить, что ключ отсутствует.
\end{enumerate}

\subsection{Открытая адресация}

Все элементы хранятся внутри массива. При коллизии ищется другая ячейка.

Линейное пробирование:

\[
h_i(k) = (h(k)+i)\bmod m.
\]

Квадратичное пробирование:

\[
h_i(k) = (h(k)+c_1i+c_2i^2)\bmod m.
\]

Двойное хеширование:

\[
h_i(k) = (h_1(k)+i h_2(k))\bmod m.
\]

\subsection{Коэффициент заполнения}

\begin{definition}
\textbf{Коэффициент заполнения} хеш-таблицы:

\[
\alpha=\frac{n}{m},
\]

где $n$ --- число элементов, $m$ --- число ячеек.
\end{definition}

При слишком большом $\alpha$ операции становятся медленнее. Обычно при превышении порога выполняют расширение таблицы и перехеширование.

\section{Связь этапов анализа}

Рассмотрим пример:

\begin{lstlisting}
{
    int x;
    x = 2 + 3 * 4;
}
\end{lstlisting}

\subsection{Лексический анализ}

Получаются токены:

\[
\{,\quad int,\quad id(x),\quad ;,\quad id(x),\quad =,\quad num(2),\quad +,\quad num(3),\quad *,\quad num(4),\quad ;,\quad \}.
\]

\subsection{Синтаксический анализ}

Строится структура:

\[
\begin{forest}
[Block
  [Decl [int] [x]]
  [Assign
    [x]
    [+
      [2]
      [*
        [3]
        [4]
      ]
    ]
  ]
]
\end{forest}
\]

\subsection{Семантический анализ}

\begin{enumerate}
    \item При входе в блок создается новая область видимости.
    \item Объявление \texttt{int x;} добавляет \texttt{x} в таблицу символов.
    \item В операторе присваивания имя \texttt{x} успешно находится.
    \item Проверяется тип правой части:
    \[
    2+3*4 : int.
    \]
    \item Проверяется совместимость:
    \[
    int := int.
    \]
    \item При выходе из блока локальная таблица удаляется.
\end{enumerate}

\section{Итоговая схема фронтенда компилятора}

\[
\begin{array}{c}
\text{Исходный текст} \\
\downarrow \\
\text{Сканер: регулярные выражения, конечные автоматы} \\
\downarrow \\
\text{Поток токенов} \\
\downarrow \\
\text{Парсер: КС-грамматика, LL(1) или LR(1)} \\
\downarrow \\
\text{Дерево разбора или AST} \\
\downarrow \\
\text{Семантический анализ: атрибуты, таблица символов} \\
\downarrow \\
\text{Промежуточное представление}
\end{array}
\]

\section{Краткий список ключевых понятий}

\begin{itemize}
    \item Регулярные языки описываются регулярными выражениями, регулярными грамматиками и конечными автоматами.
    \item Лексический анализ удобно реализуется с помощью ДКА.
    \item Контекстно-свободные грамматики описывают вложенную синтаксическую структуру программ.
    \item LL(1)-анализ строит левосторонний вывод сверху вниз.
    \item LR(1)-анализ строит правосторонний вывод в обратном порядке снизу вверх.
    \item FIRST и FOLLOW нужны для построения LL-таблиц и анализа выбора продукций.
    \item LR(1)-пункты, CLOSURE и GOTO используются для построения LR-таблиц.
    \item Атрибутные грамматики связывают синтаксис с семантическими вычислениями.
    \item AST хранит существенную структуру программы, в отличие от полного дерева разбора.
    \item Таблица символов обеспечивает работу с областями видимости, типами и объявлениями.
    \item Хеш-таблицы часто используются как эффективная основа таблицы символов.
\end{itemize}

\end{document}
